\documentclass[../../main.tex]{subfiles}
\begin{document}
\begin{oframed}
  \subfile{problem}
\end{oframed}

{\bf [解答]}

\paragraph{(1)}

% グラフの概要図
\begin{figure}[htb]
  \centering
  \begin{tikzpicture}[x=2.2cm, y=1.0cm, >=stealth, every node/.style={font=\small}]

    % ============================================================
    % 1. 定数・幾何パラメータ設定
    % 条件を満たす厳密解 f(t) = (64/15) * t^(-5)
    % x = 1.0, 2x = 2.0 のとき S(1) = 1
    % ============================================================
    \def\xVal{1.0}
    \def\twoX{2.0}
    \pgfmathsetmacro{\cFactor}{64.0 / 15.0}              % ≈ 4.267
    \pgfmathsetmacro{\fxVal}{\cFactor / pow(\xVal, 5)}   % ≈ 4.267
    \pgfmathsetmacro{\fTwoX}{\cFactor / pow(\twoX, 5)}  % ≈ 0.133

    \coordinate (O)  at (0, 0);
    \coordinate (Px) at (\xVal, 0);
    \coordinate (Qx) at (\twoX, 0);
    \coordinate (P)  at (\xVal, \fxVal);
    \coordinate (Q)  at (\twoX, \fTwoX);

    % ============================================================
    % 2. 面積 S(x) の斜線塗りつぶし (x から 2x までの曲線下)
    % ============================================================
    \fill[pattern=north east lines, pattern color=blue!70, opacity=0.75]
    (\xVal, 0) -- (P)
    -- plot[domain=\xVal:\twoX, samples=60, smooth, variable=\t] ({\t}, {\cFactor / pow(\t, 5)})
    -- (\twoX, 0) -- cycle;

    % ============================================================
    % 3. 動径 OP, OQ で挟まれる三角形 (3S(x) の幾何的等式: 3S = S + △OPP_x - △OQQ_x)
    % ============================================================
    \fill[pattern=north west lines, pattern color=red!60, opacity=0.6] (O) -- (P) -- (Px) -- cycle;
    \fill[pattern=crosshatch dots, pattern color=teal!70, opacity=0.6] (O) -- (Q) -- (Qx) -- cycle;

    % ============================================================
    % 4. 座標軸
    % ============================================================
    \draw[->, thick, gray!80!black] (-0.2, 0) -- (\twoX + 0.6, 0) node[right] {$t$};
    \draw[->, thick, gray!80!black] (0, -0.3) -- (0, \fxVal + 0.7) node[above] {$y$};
    \node[below left=2pt] at (O) {$O$};

    % ============================================================
    % 5. 曲線 y = f(t) (単調減少)
    % ============================================================
    \draw[line width=1.6pt, blue!85!black, domain=0.92:2.35, samples=70, smooth, variable=\t]
    plot ({\t}, {\cFactor / pow(\t, 5)});
    \node[blue!90!black, font=\footnotesize\bfseries, above right] at (1.5, 0.9)
    {$y = f(t)$};

    % ============================================================
    % 6. 動径 OP, OQ および垂直線分
    % ============================================================
    % 動径 OP, OQ
    \draw[line width=1.2pt, red!85!black] (O) -- (P);
    \draw[line width=1.2pt, teal!85!black] (O) -- (Q);

    % 垂線 (P -> Px, Q -> Qx)
    \draw[densely dotted, thick, gray!70!black] (P) -- (Px);
    \draw[densely dotted, thick, gray!70!black] (Q) -- (Qx);

    % y 軸への投影破線
    \draw[densely dotted, thick, gray!70!black] (P) -- (0, \fxVal) node[left=2pt, font=\scriptsize] {$f(x)$};
    \draw[densely dotted, thick, gray!70!black] (Q) -- (0, \fTwoX) node[left=2pt, font=\scriptsize] {$f(2x)$};

    % ============================================================
    % 7. 点プロットと目盛りラベル
    % ============================================================
    \node[circle, fill, inner sep=1.3pt] at (P) {}; \node[above right=1pt, font=\footnotesize\bfseries] at (P) {$P(x, f(x))$};
    \node[circle, fill, inner sep=1.3pt] at (Q) {}; \node[above right=1pt, font=\footnotesize\bfseries] at (Q) {$Q(2x, f(2x))$};

    \node[below=2pt, font=\small\bfseries] at (Px) {$x$};
    \node[below=2pt, font=\small\bfseries] at (Qx) {$2x$};

    % 面積 S(x) のラベル
    \node[blue!95!black, font=\large\bfseries] at (1.35, 1.4) {$S(x)$};

  \end{tikzpicture}
  \caption{単調減少関数 $y = f(t)$，面積 $S(x)$ および各領域}
  \label{fig:decreasing_function_and_area_Sx}
\end{figure}

題意から，$f(x)$は$0<x$で単調減少な正の値をとる関数である．
また$S(a)$は
\begin{align}
  S(x)=\int_x^{2x}f(t)dt \label{eq:1}
\end{align}
である．

次に，$3S(a)$に関する面積の条件より
\begin{align}
  3S(x)
  &=\int_x^{2x}f(t)dt+\dfrac{1}{2}xf(x)-\dfrac{1}{2}2xf(2x) \nonumber\\
  &=S(x)+\dfrac{1}{2}xf(x)-xf(2x) \nonumber\\
  \therefore
  S(x)
  &=\dfrac{1}{4}x(f(x)-2f(2x)) \label{eq:2}
\end{align}
である．

\cref{eq:1}の両辺を$x$で微分すると
\begin{align}
  S'(x)=2f(2x)-f(x) \label{eq:3}
\end{align}
だから，\cref{eq:2}と合わせて$f(x)$を消去すると，$y=S(x)$として
\begin{align}
  4y=-x\dfrac{dy}{dx}
\end{align}
である．$0<x$で$0<f(x)$だから\cref{eq:1}より$y>0$となるので，変形して
\begin{align}
  \dfrac{-4dx}{x}=\dfrac{dy}{y}
\end{align}
である．積分して題意の条件$S(1)=1$より，
\begin{align}
  y=S(x)=x^{-4}
\end{align}
である．\cref{eq:2}から
\begin{align}
  f(x)-2f(2x) = \dfrac{4S(x)}{x} = 4x^{-5}
\end{align}
である．以上まとめて
\begin{align}
  \begin{cases}
    S(x)=x^{-4} \\
    f(x)-2f(2x) = 4x^{-5}
  \end{cases} \label{eq:4}
\end{align}
である．$\cdots$(答)

\paragraph{(2)}

$b_n(x)=2^nf(2^nx)$とおくと，初期条件$b_0=f(x)$である．
\cref{eq:4}の第二式に$x$に$2^nx$を代入して，
\begin{align}
  & 2f(2^{n+1}x)=f(2^nx)-4(2^nx)^{-5} \\
  & 2^{n+1}f(2^{n+1}x)=2^nf(2^nx)-2^{n+2}(2^nx)^{-5} \\
  \therefore
  & b_{n+1}(x)=b_n(x)-2^{2-4n}x^{-5} \label{eq:5}
\end{align}
である．$b_n(x)$が階差数列だから，$n\ge1$のとき
\begin{align}
  & b_n(x)=b_0(x)-\sum_{k=0}^{n-1}2^{2-4k}x^{-5} \\
  & b_n(x)=f(x)-4x^{-5}\dfrac{1-2^{-4n}}{1-2^{-4}} \\
  \therefore
  & \lim_{n\to\infty}b_n = f(x)-\dfrac{64}{15}x^{-5}=a(x) \label{eq:6}
\end{align}
である．

従って求める積分値は
\begin{align}
  \int_{x}^{2x}a(t)dt
  &=S(x)+\left[\dfrac{16}{15}t^{-4}\right]_x^{2x} \\
  &=x^{-4}+\dfrac{16}{15}\left(\dfrac{1}{16}-1\right)x^{-4} \\
  &=0
\end{align}
である．

\paragraph{(3)}

$f(x)>0$だから$a(x) \ge 0$である．これと(2)の積分計算から被積分関数$a(x)$が恒等的に$0$である．
従って\cref{eq:6}から
\begin{align}
  f(x)-\dfrac{64}{15}x^{-5} = 0 \quad \therefore
  f(x) = \dfrac{64}{15}x^{-5}
\end{align}
である．
これは$x>0$で定義された連続な単調減少関数であり，$S(x)$に関する面積条件も満たすから十分である．
以上から，
\begin{align}
  f(x)=\dfrac{64}{15}x^{-5}
\end{align}
である．$\cdots$(答)

\end{document}
